% ===================================================================
%  common/preamble.tex — 全部课程共用导言区（论文风格版）
%  引擎：XeLaTeX ｜ 正文：宋体 + TeX Gyre Termes ｜ 数学：Termes Math
%  规则：未经同意不在此处增删宏包
% ===================================================================

% ---------- 版心 ----------
\usepackage[a4paper,left=2.6cm,right=2.6cm,top=2.6cm,bottom=2.4cm]{geometry}

% ---------- 字体 ----------
\usepackage{unicode-math}
\setmainfont{TeX Gyre Termes}
\setmathfont{TeX Gyre Termes Math}
\setCJKmainfont{SimSun}[BoldFont=SimHei,ItalicFont=KaiTi]
\setCJKsansfont{SimHei}
\setCJKmonofont{FangSong}
% 只用西文标点：禁用 xeCJK 的全角标点替换与挤压
\punctstyle{plain}
\linespread{1.35}
\setlength{\parskip}{0.2em}
\setlength{\parindent}{2em}   % 首行缩进两个汉字

% ---------- 数学 ----------
\usepackage{amsmath}
\usepackage{graphicx}
\usepackage{booktabs}
\usepackage{array}

% ---------- 标题格式：居中、黑体 ----------
\usepackage{indentfirst}
\usepackage{titlesec}
\titleformat{\section}{\centering\heiti\bfseries\large}{\thesection}{0.6em}{}
\titleformat{\subsection}{\heiti\bfseries\normalsize}{\thesubsection}{0.6em}{}
\titlespacing*{\section}{0pt}{3ex plus .6ex minus .3ex}{1.6ex}
\titlespacing*{\subsection}{0pt}{2.8ex plus .5ex minus .2ex}{1.1ex}

% ---------- 页眉页脚 ----------
\usepackage{fancyhdr}
\pagestyle{fancy}
\fancyhf{}
\fancyhead[C]{\small\kaishu\leftmark}
\fancyfoot[C]{\small\thepage}
\renewcommand{\headrulewidth}{0.4pt}
\fancypagestyle{plain}{\fancyhf{}\fancyfoot[C]{\small\thepage}\renewcommand{\headrulewidth}{0pt}}
\renewcommand{\sectionmark}[1]{}   % 页眉固定为课程名

% ---------- 列表：论文式 (1) (2) ----------
\usepackage{enumitem}
\setlist[enumerate,1]{label=(\arabic*),leftmargin=2.2em,itemsep=0.2em,topsep=0.3em}
\setlist[enumerate,2]{label=(\roman*),leftmargin=1.8em,itemsep=0.1em,topsep=0.2em}
\setlist[itemize,1]{leftmargin=2em}

% ---------- 定理环境 ----------
\usepackage{amsthm}
\newtheoremstyle{cnplain}
  {0.9em}   % 上方间距
  {0.9em}   % 下方间距
  {\normalfont}  % 正文字体
  {}        % 缩进
  {\heiti\bfseries}  % 标题字体
  {}        % 标题后标点
  {0.7em}   % 标题后间距
  {\thmname{#1}\thmnumber{ #2}\thmnote{\normalfont\ (#3)}}

\theoremstyle{cnplain}
\newtheorem{theorem}{定理}[section]
\newtheorem{lemma}[theorem]{引理}
\newtheorem{proposition}[theorem]{命题}
\newtheorem{corollary}[theorem]{推论}
\newtheorem{definition}[theorem]{定义}
\newtheorem{example}[theorem]{例}
\newtheorem{remark}[theorem]{注}
\newtheorem{exercise}[theorem]{习题}

\renewcommand{\proofname}{\heiti 证明}
\renewcommand{\qedsymbol}{\rule{0.55em}{0.55em}}

% ---------- 超链接 ----------
\usepackage{hyperref}
\hypersetup{
  colorlinks=true,
  linkcolor=[rgb]{0.10,0.20,0.55},
  citecolor=[rgb]{0.10,0.35,0.20},
  urlcolor=[rgb]{0.10,0.20,0.55},
  bookmarksnumbered=true,
}
\setlength{\emergencystretch}{3em}
\makeatletter
\def\@doendpe{\@endpetrue}   % 列表结束后不取消首行缩进
\makeatother

